% !TeX root = main.tex
Multi-fingered, contact-rich manipulation over long execution horizons remains
one of the open frontiers of real-world robotics~\cite{siciliano2016springer}.
Tasks such as threading a nut onto a bolt are effortless for humans yet expose
the full difficulty of autonomous dexterity: the robot must regulate motion
through sustained contact under tight geometric tolerances, uncertain friction,
and only partially observed interaction forces. Small pose errors cause
slipping, jamming, or cross-threading, while excessive force damages the
workpiece. These tightly coupled perception, control, and contact constraints
make controlled threading one of the most demanding contact-rich manipulation
problems, and its difficulty is compounded when a dexterous hand, rather than a
parallel gripper, must coordinate many degrees of freedom to sustain
engagement.

The difficulty grows sharply over long horizons. Because the reachable rotation
of a wrist or hand is bounded, a single continuous turn cannot fully seat a nut;
the operation instead requires a \emph{cyclic} sequence of partial rotations,
timely release, collision-free return, and re-engagement, repeated until the nut
is fully threaded. This cyclic structure is what turns a short-horizon contact
skill into a long-horizon task: the terminal state of each cycle is the initial
condition of the next, so a small deviation in one rotation perturbs contact in
the following one. Under closed-loop autoregressive execution such deviations
accumulate, driving the policy toward out-of-distribution states from which
recovery is unlikely---the compounding-error, or covariate-shift, failure mode
that is well understood in imitation learning~\cite{ross2011reduction}. Purely
scripted trajectories cannot react to these deviations, while a monolithic
learned policy must simultaneously discover fine contact behavior and the
long-horizon phase transitions that stitch cycles together.

Sim-to-real reinforcement learning with privileged information has become the
dominant recipe for short-horizon contact-rich skills. High-fidelity contact
simulators built on signed-distance-field (SDF) collision
representations~\cite{narang2022factory} enable stable, fast simulation of tight
assembly, and force-aware policies trained under dynamics randomization transfer
robustly to hardware for tasks including insertion, gear meshing, and single
nut threading~\cite{noseworthy2025forge}. However, these successes are largely
confined to short contact episodes. When the same recipe is scaled to complex
dexterous assets and long horizons, two failure modes emerge. First, residual
gaps in contact and friction modeling let a reinforcement-learning agent
\emph{exploit} simulator inaccuracies rather than learn behavior that survives
the sim-to-real gap---a pathology observed precisely in dexterous nut
screwing~\cite{dexscrew}. Second, the trajectory drift described above turns
otherwise reliable per-cycle policies into unreliable long-horizon ones.

The conventional remedy for long horizons is \emph{policy chaining}:
decomposing the task into discrete sub-skills and learning to transition between
them. Sequential Dexterity~\cite{chen2023sequential}, the representative
approach, chains several reinforcement-learning sub-policies using a learned
transition-feasibility function. This modular design is expressive, but it
carries structural costs that are particularly acute for a cyclic task: it
requires training and maintaining multiple networks, hand-designing or learning
the transition boundaries between them, and absorbing the boundary failures and
execution latency that arise whenever control is handed from one sub-policy to
the next. For a task whose sub-skills are nearly identical and must repeat many
times, replicating this discrete machinery per cycle is both wasteful and a new
source of compounding boundary error.

We argue that the cyclic structure of threading is better captured by a single
\emph{continuous} generative policy than by a chain of discrete ones. We
introduce \methodname, a framework that distills a privileged simulation
teacher into a unified vision-and-force flow-matching student for long-horizon
cyclic threading (\figref{fig:teaser}). We first train a reactive teacher in a
simplified simulation that preserves the SDF-based contact interactions
essential to threading while removing irrelevant complexity, giving the teacher
access to privileged state, contact forces, and exact geometry. Because a single
reachable rotation is bounded, we generate long-horizon demonstrations by rolling
out this \emph{single} reactive teacher inside a state machine: the policy
threads until a rotation threshold is reached, a guided open-and-return motion
then resets the hand to its initial opened configuration, and the same policy is
re-invoked to thread the next increment in the same environment. Repeating this
cycle produces arbitrarily long demonstrations without ever switching between
distinct learned sub-policies. We then distill these long-horizon rollouts into a
single student policy that replaces privileged state with a DP3 point-cloud
vision head~\cite{ze2024dp3} and force feedback, and that models the full cyclic
trajectory distribution with flow-matching imitation
learning~\cite{lipman2023flow}. Flow matching is well suited to this setting:
it regresses a straight, deterministic velocity field between noise and expert
actions, yielding a multi-modal yet efficient trajectory generator that captures
long-horizon behavior in a single network. Crucially, the deployed student is
\emph{one} continuous policy---it reproduces the entire multi-cycle chain
without switching between discrete sub-networks at runtime, so it avoids the
inter-policy boundary failures and latency of explicit chaining while remaining
reactive to contact.

The main contributions of this work are:
\begin{itemize}
  \item \textbf{System and task.} A hardware platform and benchmark for
    long-horizon cyclic nut threading on a 7-DoF arm equipped with a
    multi-fingered dexterous hand, exercising force-regulated contact over many
    rotation--release--return--re-engage cycles.
  \item \textbf{Privileged flow-matching distillation.} A teacher-to-student
    pipeline that trains a privileged reactive teacher in a simplified
    SDF-contact simulation and distills its long-horizon behavior into a single
    flow-matching policy, rather than a chain of learned sub-policies.
  \item \textbf{Cross-modal sim-to-real integration.} A student policy that
    substitutes a DP3 visual-latent representation and force feedback for the
    teacher's privileged simulation state, transferring contact-rich behavior to
    the real world without runtime policy switching.
  \item \textbf{Empirical validation.} An evaluation across M24, M30, and M36
    nuts with hexagonal, octagonal, rectangular, and round geometries, reporting
    maximum cycle count, threading tracking error, and success rate to quantify
    robustness and generalization.
\end{itemize}
